% ========================================================
% BÀI TẬP TOÁN 5 - KẾT NỐI TRI THỨC VỚI CUỘC SỐNG
% Bài 61: Thực hành tính toán và ước lượng vận tốc, quãng đường, thời gian
% Định dạng: BTTeX Part III - Trắc nghiệm điền số ngắn
% ========================================================

% Câu 1
\begin{ex}
  Hai người xuất phát cùng lúc từ $A$ và $B$ cách nhau $120\text{ km}$, đi ngược chiều nhau. Vận tốc người thứ nhất là $30\text{ km/h}$, người thứ hai là $10\text{ km/h}$. Sau mấy giờ hai người gặp nhau?
  \shortans{3}
  \loigiai{
    Tổng vận tốc: $30 + 10 = 40\text{ km/h}$. Thời gian gặp: $120 : 40 = 3\text{ giờ}$.
  }
\end{ex}

% Câu 2
\begin{ex}
  Quãng đường $AB$ dài $180\text{ km}$. Hai ô tô đi ngược chiều nhau: xe từ $A$ đi $50\text{ km/h}$, xe từ $B$ đi $40\text{ km/h}$. Sau mấy giờ hai xe gặp nhau?
  \shortans{2}
  \loigiai{
    Tổng vận tốc: $50 + 40 = 90\text{ km/h}$. Thời gian gặp: $180 : 90 = 2\text{ giờ}$.
  }
\end{ex}

% Câu 3
\begin{ex}
  Ở câu trên, điểm gặp nhau cách $A$ bao nhiêu ki-lô-mét?
  \shortans{100}
  \loigiai{
    $50 \times 2 = 100\text{ km}$.
  }
\end{ex}

% Câu 4
\begin{ex}
  Hai người đi xe máy khởi hành cùng lúc từ hai địa điểm cách nhau $150\text{ km}$ đi ngược chiều và gặp nhau sau $2$ giờ. Tổng vận tốc của hai người là bao nhiêu km/h?
  \shortans{75}
  \loigiai{
    $150 : 2 = 75\text{ km/h}$.
  }
\end{ex}

% Câu 5
\begin{ex}
  Một người đi xe đạp khởi hành từ $A$ với vận tốc $12\text{ km/h}$. Sau $2$ giờ, một xe máy cũng đi từ $A$ đuổi theo với vận tốc $36\text{ km/h}$. Sau mấy giờ xe máy đuổi kịp xe đạp?
  \shortans{1}
  \loigiai{
    Khoảng cách lúc xe máy xuất phát: $12 \times 2 = 24\text{ km}$. Hiệu vận tốc: $36 - 12 = 24\text{ km/h}$. Thời gian đuổi kịp: $24 : 24 = 1\text{ giờ}$.
  }
\end{ex}

% Câu 6
\begin{ex}
  Một ca nô có vận tốc thực khi nước yên lặng là $25\text{ km/h}$, vận tốc dòng nước là $3\text{ km/h}$. Vận tốc ca nô khi xuôi dòng là bao nhiêu km/h?
  \shortans{28}
  \loigiai{
    $25 + 3 = 28\text{ km/h}$.
  }
\end{ex}

% Câu 7
\begin{ex}
  Ở câu trên, vận tốc ca nô khi ngược dòng là bao nhiêu km/h?
  \shortans{22}
  \loigiai{
    $25 - 3 = 22\text{ km/h}$.
  }
\end{ex}

% Câu 8
\begin{ex}
  Một con thuyền đi xuôi dòng quãng đường $40\text{ km}$ với vận tốc $20\text{ km/h}$. Thời gian thuyền đi hết bao nhiêu giờ?
  \shortans{2}
  \loigiai{
    $40 : 20 = 2\text{ giờ}$.
  }
\end{ex}

% Câu 9
\begin{ex}
  Hai thành phố cách nhau $240\text{ km}$. Một ô tô đi từ $A$ đến $B$ với vận tốc $60\text{ km/h}$, đồng thời một ô tô đi từ $B$ đến $A$ với vận tốc $40\text{ km/h}$. Sau bao nhiêu phút hai xe gặp nhau?
  \shortans{144}
  \loigiai{
    Tổng vận tốc: $100\text{ km/h}$. Thời gian gặp: $240 : 100 = 2{,}4\text{ giờ} = 144$ phút.
  }
\end{ex}

% Câu 10
\begin{ex}
  Hai xe đạp khởi hành cùng lúc từ cùng một điểm đi cùng chiều: xe thứ nhất đi $15\text{ km/h}$, xe thứ hai đi $11\text{ km/h}$. Sau $2$ giờ hai xe cách nhau bao nhiêu ki-lô-mét?
  \shortans{8}
  \loigiai{
    Mỗi giờ cách nhau $15 - 11 = 4\text{ km}$. Sau 2 giờ cách $4 \times 2 = 8\text{ km}$.
  }
\end{ex}

% Câu 11
\begin{ex}
  Một tàu thủy xuôi dòng với vận tốc $28\text{ km/h}$ và ngược dòng với vận tốc $22\text{ km/h}$. Vận tốc của dòng nước là bao nhiêu km/h?
  \shortans{3}
  \loigiai{
    $(28 - 22) : 2 = 3\text{ km/h}$.
  }
\end{ex}

% Câu 12
\begin{ex}
  Ở câu trên, vận tốc thực của tàu thủy khi nước yên lặng là bao nhiêu km/h?
  \shortans{25}
  \loigiai{
    $(28 + 22) : 2 = 25\text{ km/h}$.
  }
\end{ex}

% Câu 13
\begin{ex}
  Đoạn đường sắt dài $300\text{ m}$. Một đoàn tàu dài $200\text{ m}$ chạy qua đoạn đường đó với vận tốc $20\text{ m/s}$. Thời gian tàu chạy qua hết bao nhiêu giây?
  \shortans{25}
  \loigiai{
    Quãng đường đi: $300 + 200 = 500\text{ m}$. Thời gian: $500 : 20 = 25\text{ giây}$.
  }
\end{ex}

% Câu 14
\begin{ex}
  Một ô tô đi quãng đường $120\text{ km}$. Nửa quãng đường đầu đi với vận tốc $60\text{ km/h}$, nửa quãng đường sau đi $40\text{ km/h}$. Tổng thời gian đi hết bao nhiêu phút?
  \shortans{150}
  \loigiai{
    Thời gian đầu: $60 : 60 = 1$ giờ; thời gian sau: $60 : 40 = 1{,}5$ giờ. Tổng $2{,}5\text{ giờ} = 150$ phút.
  }
\end{ex}

% Câu 15
\begin{ex}
  Một người đi bộ $1$ vòng hồ dài $1\,200\text{ m}$ hết $15$ phút. Vận tốc đi bộ là bao nhiêu m/phút?
  \shortans{80}
  \loigiai{
    $1\,200 : 15 = 80\text{ m/phút}$.
  }
\end{ex}

% Câu 16
\begin{ex}
  Một xe tải và một xe máy cùng đi từ $A$ đến $B$. Xe máy đi $40\text{ km/h}$, xe tải đi $50\text{ km/h}$. Sau $3$ giờ xe tải đi trước xe máy bao nhiêu ki-lô-mét?
  \shortans{30}
  \loigiai{
    $(50 - 40) \times 3 = 30\text{ km}$.
  }
\end{ex}

% Câu 17
\begin{ex}
  Hai thị trấn $A$ và $B$ cách nhau $64\text{ km}$. Một người đi xe đạp từ $A$ lúc $7\text{ giờ}$ với vận tốc $16\text{ km/h}$. Người đó đến $B$ lúc mấy giờ?
  \shortans{11}
  \loigiai{
    Thời gian đi: $64 : 16 = 4\text{ giờ}$. Đến nơi lúc $7 + 4 = 11\text{ giờ}$.
  }
\end{ex}

% Câu 18
\begin{ex}
  Một máy bay dự kiến bay từ Hà Nội đến TP Hồ Chí Minh dài $1\,200\text{ km}$ với vận tốc $600\text{ km/h}$. Chuyến bay bay hết bao nhiêu giờ?
  \shortans{2}
  \loigiai{
    $1\,200 : 600 = 2\text{ giờ}$.
  }
\end{ex}

% Câu 19
\begin{ex}
  Ước lượng vận tốc người đi bộ bình thường khoảng $4\text{ km/h}$. Để đi $12\text{ km}$ người đó cần bao nhiêu giờ?
  \shortans{3}
  \loigiai{
    $12 : 4 = 3\text{ giờ}$.
  }
\end{ex}

% Câu 20
\begin{ex}
  Một cano đi từ bến $A$ đến bến $B$ cách nhau $36\text{ km}$ xuôi dòng mất $1{,}5$ giờ. Vận tốc dòng nước là $4\text{ km/h}$. Vận tốc thực của ca nô là bao nhiêu km/h?
  \shortans{20}
  \loigiai{
    Vận tốc xuôi: $36 : 1{,}5 = 24\text{ km/h}$. Vận tốc thực: $24 - 4 = 20\text{ km/h}$.
  }
\end{ex}

% Câu 21
\begin{ex}
  Ở câu trên, vận tốc ngược dòng của ca nô là bao nhiêu km/h?
  \shortans{16}
  \loigiai{
    $20 - 4 = 16\text{ km/h}$.
  }
\end{ex}

% Câu 22
\begin{ex}
  Hai người cùng đi từ hai đầu một cây cầu dài $450\text{ m}$ ngược chiều nhau. Người thứ nhất đi $2\text{ m/s}$, người thứ hai đi $1\text{ m/s}$. Sau bao nhiêu giây hai người gặp nhau?
  \shortans{150}
  \loigiai{
    Tổng vận tốc: $3\text{ m/s}$. Thời gian: $450 : 3 = 150\text{ giây}$.
  }
\end{ex}

% Câu 23
\begin{ex}
  Một ô tô đi từ Hà Nội về Nam Định với vận tốc $60\text{ km/h}$ hết $1\text{ giờ } 30\text{ phút}$. Quãng đường dài bao nhiêu ki-lô-mét?
  \shortans{90}
  \loigiai{
    $60 \times 1{,}5 = 90\text{ km}$.
  }
\end{ex}

% Câu 24
\begin{ex}
  Bác An chạy bộ quanh sân vận động có chu vi $400\text{ m}$. Bác chạy $5$ vòng hết $10$ phút. Vận tốc của bác An là bao nhiêu m/phút?
  \shortans{200}
  \loigiai{
    Quãng đường: $400 \times 5 = 2\,000\text{ m}$. Vận tốc: $2\,000 : 10 = 200\text{ m/phút}$.
  }
\end{ex}

% Câu 25
\begin{ex}
  Một con thỏ chạy với vận tốc $15\text{ m/s}$. Để chạy $300\text{ m}$ thỏ cần bao nhiêu giây?
  \shortans{20}
  \loigiai{
    $300 : 15 = 20\text{ giây}$.
  }
\end{ex}
